\documentclass[11pt,a4paper]{article}

\usepackage[margin=1in]{geometry}
\usepackage{amsmath,amssymb}
\usepackage{booktabs}
\usepackage{tabularx}
\usepackage{hyperref}
\usepackage{xcolor}
\usepackage[most]{tcolorbox}
\usepackage[htt]{hyphenat} % lets \texttt names break at underscores
\makeatletter\g@addto@macro\@verbatim\small\makeatother % code blocks fit the text width
\hypersetup{colorlinks=true,linkcolor=blue!50!black,urlcolor=blue!50!black}

\newcommand{\Mc}{\mathcal{M}}
\newcommand{\lt}{\tilde{\Lambda}}
\newcommand{\dlt}{\delta\tilde{\Lambda}}
\newcommand{\dd}{\mathrm{d}}
\newcommand{\code}[1]{\texttt{#1}}

% A boxed note for bugs found and fixed in the code, placed next to the
% derivation or mechanism they concern.
\newtcolorbox{bugnote}[1]{colback=red!3!white,colframe=red!45!black,
  fonttitle=\bfseries,title={Bug (fixed): #1},breakable}
% A boxed note for known limitations that are not fixed.
\newtcolorbox{caveat}[1]{colback=yellow!4!white,colframe=yellow!45!black,
  coltitle=black,fonttitle=\bfseries,title={Caveat: #1},breakable}

\title{TOV-informed joint priors in \code{bilby.gw.prior}:\\
$(m, \Lambda)$ and $(\Mc,\, q,\, \lt,\, \dlt)$}
\author{Derivation and implementation notes}
\date{\today}

\begin{document}
\sloppy
\maketitle

\begin{abstract}
\noindent
The equation of state (EOS) correlates a neutron star's mass with its tidal
deformability $\Lambda$. These notes document three priors in
\code{bilby/gw/prior.py} that carry that correlation from a gridded KDE of TOV
mass--$\Lambda$ sequences into a bilby run. \code{TOVJointDist} and
\code{TOVJointPrior} put it on a single $(m, \Lambda)$ pair, one instance per
star. \code{TOVMassLambdaTildeJointDist} and \code{TOVMassLambdaTildeJointPrior}
push two such pairs through the exact change of variables to the usual BNS
sampling parameters $(\Mc, q, \lt, \dlt)$. \code{TOVConditionalLambdaPrior}
puts only the conditional $p(\Lambda \mid m)$ on each star's deformability and
leaves the masses to whatever prior the run already uses. For each, we derive
the density evaluated by \code{ln\_prob} and the unit-cube transform used by
\code{rescale}, show that the two agree, and describe how the classes plug into
bilby's prior machinery. Bugs found and fixed along the way are recorded
where they occur.
\end{abstract}

\tableofcontents

%======================================================================
\section{Overview}
%======================================================================

\begin{table}[h]
\centering
\begin{tabularx}{\textwidth}{lX}
\toprule
Class / function & Role \\
\midrule
\code{TOVJointDist(grid\_file, names, bounds=None)} &
  Joint distribution of one star's $(m, \Lambda)$; owns the grid and all the
  tables derived from it (Sections~\ref{sec:grid}--\ref{sec:single}). \\
\code{TOVJointPrior(dist, name)} &
  Per-parameter \code{JointPrior} wrapper; two per star. \\
\code{TOVMassLambdaTildeJointDist(grid\_file, names=None, bounds=None)} &
  Joint distribution of $(\Mc, q, \lt, \dlt)$ for a binary; holds a
  \code{TOVJointDist} internally as \code{\_grid} and reuses its tables
  (Section~\ref{sec:binary}). \\
\code{TOVMassLambdaTildeJointPrior(dist, name)} &
  Per-parameter \code{JointPrior} wrapper; four per binary. \\
\code{TOVConditionalLambdaPrior(grid\_file, component, ...)} &
  Conditional prior $p(\Lambda \mid m)$ on one star's deformability, for use
  with an independent mass prior; one per star (Section~\ref{sec:conditional}). \\
\code{\_lambda\_tilde\_coefficients(eta)} &
  The coefficients $c_1,\dots,c_4$ of the tidal map (Eq.~\ref{eq:coeffs}). \\
\code{\_ln\_abs\_det\_lambda\_tilde\_jacobian(eta)} &
  $\ln\lvert\det A(\eta)\rvert$ of the tidal map (Eq.~\ref{eq:tidaljac}). \\
\bottomrule
\end{tabularx}
\caption{The public surface. Both \code{Dist} classes derive from
\code{bilby.core.prior.BaseJointPriorDist}, and both joint \code{Prior} classes
derive from \code{bilby.core.prior.JointPrior}.
\code{TOVConditionalLambdaPrior} derives from
\code{bilby.core.prior.ConditionalBasePrior} instead, because it conditions on
parameters it does not own.}
\end{table}

\noindent
Each \code{Dist} implements the three methods bilby expects of a joint
distribution: \code{\_ln\_prob} (density), \code{\_rescale} (unit-cube
transform, used by nested samplers) and \code{\_sample} (random draws). A
sampler only ever reaches them through the per-parameter \code{Prior} objects
and a \code{PriorDict}; Section~\ref{sec:machinery} covers that layer.

%======================================================================
\section{The input grid}
\label{sec:grid}
%======================================================================

\subsection{File format}

The EOS information enters through a single file,
\code{prob\_data/tov\_mass\_lambda\_joint\_kde\_grid.npz}, loaded in
\code{TOVJointDist.\_\_init\_\_}. It holds a two-dimensional KDE of TOV
sequences over a population of equations of state:
\begin{equation}
  \code{log\_pdf}[i,j] \;=\; \ln p\big(m_i,\, \ell_j\big),
  \qquad \ell \equiv \log_{10}\Lambda ,
\end{equation}
on grids \code{mass\_edges} ($N=200$ points, $m \in [1.000, 3.476]\,M_\odot$)
and \code{log10\_lambda\_edges} ($M=200$ points, $\ell \in [0.568, 4.738]$). The
grid is normalised, $\int p(m,\ell)\,\dd m\,\dd\ell = 1$. It is fitted offline
because refitting the KDE to the full population per job, let alone per
likelihood call, is prohibitive. Empty regions of the KDE are floored at
$\ln p = \ln 10^{-300} = -690.8$ rather than $-\infty$.

\subsection{Tables derived at construction}
\label{sec:tables}

\code{TOVJointDist.\_\_init\_\_} builds three objects from the grid, all by
trapezoidal quadrature (\code{scipy.integrate.trapezoid} and
\code{cumulative\_trapezoid}).

\paragraph{The density interpolant.}
\code{self.\_log\_pdf\_interp} is a SciPy \code{RegularGridInterpolator}
over $(m, \ell)$, linear in $\ln p$ (so geometric in $p$), with
\code{fill\_value=-inf} off the grid. It is used only by \code{ln\_prob}.

\paragraph{The marginal mass CDF.}
\begin{equation}
  p(m) = \int p(m,\ell)\,\dd\ell,
  \qquad
  F(m) = \frac{\int_{m_{\min}}^{m} p(m')\,\dd m'}{\int_{m_{\min}}^{m_{\max}} p(m')\,\dd m'} ,
  \label{eq:massmarg}
\end{equation}
tabulated at the mass grid points as \code{self.\_mass\_cdf} and explicitly
renormalised so that $F(m_{\max}) = 1$. It is exposed as
\code{TOVJointDist.mass\_cdf(mass)} (forward, linear interpolation) and
\code{TOVJointDist.mass\_from\_cdf(u)} (inverse, linear interpolation of the
same table with the axes swapped).

\paragraph{The conditional CDF of $\ell$ given $m$.}
\begin{equation}
  p(\ell \mid m_i) = \frac{p(m_i,\ell)}{p(m_i)},
  \qquad
  G_i(\ell) = \int_{\ell_{\min}}^{\ell} p(\ell' \mid m_i)\,\dd\ell' ,
  \label{eq:condcdf}
\end{equation}
one row per mass grid point, stored as \code{self.\_cond\_cdf} with shape
$(N, M)$ and each row renormalised to end at $1$. Rows whose marginal
$p(m_i)$ vanishes are guarded against division by zero.

\subsection{From $\log_{10}\Lambda$ to $\Lambda$}
\label{sec:log10jac}

The grid is a density in $\ell = \log_{10}\Lambda$; every density these classes
return is in $\Lambda$. With $\dd\ell/\dd\Lambda = 1/(\Lambda \ln 10)$,
\begin{equation}
  p(m, \Lambda) \;=\; \frac{p(m, \ell)}{\Lambda \ln 10},
  \qquad\text{i.e.}\qquad
  \ln p(m,\Lambda) = \ln p(m,\ell) - \ln \Lambda - \ln(\ln 10).
  \label{eq:log10jac}
\end{equation}
The last term is $\ln(\ln 10) = 0.8340$. Equation~\eqref{eq:log10jac} is
implemented once, in \code{TOVJointDist.ln\_pdf\_mass\_lambda(mass, lam)}, which
both \code{Dist} classes call.

\begin{bugnote}{$\ln 10$ in place of $\ln(\ln 10)$}
The original \code{TOVJointDist.\_ln\_prob} computed
\begin{center}
\code{log\_pdf\_val - np.log(lam) - np.log(10.0)}
\end{center}
that is, $\ln 10 = 2.3026$ where Eq.~\eqref{eq:log10jac} requires
$\ln(\ln 10) = 0.8340$. The log-density was low by the constant
\begin{equation*}
  \ln 10 - \ln(\ln 10) = \ln\!\left( \frac{10}{\ln 10} \right) = 1.4686
\end{equation*}
per $(m, \Lambda)$ pair, a factor of $4.34$ in density. Being constant, it
did not bias posteriors, but it shifted log evidences and left the density
unnormalised. It was found by the box estimator of Section~\ref{sec:verify}:
the four-parameter prior, which contains two pairs, came out at
$18.8 \approx 4.34^2$ instead of $1$. The fix moved the computation into
\code{ln\_pdf\_mass\_lambda} so that both classes share the corrected line.
\end{bugnote}

%======================================================================
\section{Single star: \code{TOVJointDist}}
\label{sec:single}
%======================================================================

\subsection{Density}

For parameters named \code{[mass\_name, lambda\_name]}, e.g.
\code{["mass\_1", "lambda\_1"]}, the prior is the grid density itself,
transformed to $\Lambda$:
\begin{equation}
  \ln p(m, \Lambda) = \ln p(m, \ell)\big|_{\text{interp}} - \ln\Lambda - \ln(\ln 10),
  \label{eq:single}
\end{equation}
evaluated by \code{TOVJointDist.\_ln\_prob}, which forwards to
\code{ln\_pdf\_mass\_lambda} and then sets out-of-bounds samples to $-\infty$.
The default \code{bounds} are the grid's own extent,
$m \in [m_{\min}, m_{\max}]$ and $\Lambda \in [10^{\ell_{\min}}, 10^{\ell_{\max}}]$,
so the support is a box and coincides with where the interpolant is defined.

\subsection{Prior transform}
\label{sec:single_rescale}

\code{TOVJointDist.\_rescale} maps $u = (u_0, u_1) \in [0,1]^2$ to $(m, \Lambda)$
by sequential inverse-CDF sampling, mass first and then $\ell$ given that mass:
\begin{equation}
  m = F^{-1}(u_0),
  \qquad
  \ell = G^{-1}(u_1 \mid m),
  \qquad
  \Lambda = 10^{\ell} ,
  \label{eq:single_rescale}
\end{equation}
which samples $p(m)\,p(\ell \mid m) = p(m, \ell)$, as required. The first step is
\code{mass\_from\_cdf}; the second is
\code{log10\_lambda\_from\_cdf(u, mass)}. Since $G$ is tabulated only at the mass
grid points, the latter blends the two neighbouring rows linearly in mass
before inverting:
\begin{equation}
  G(\ell \mid m) = (1-w)\, G_{i}(\ell) + w\, G_{i+1}(\ell),
  \qquad
  w = \frac{m - m_i}{m_{i+1} - m_i}, \quad m_i \le m < m_{i+1}.
  \label{eq:rowblend}
\end{equation}
The blended row is still a monotone CDF, so the inversion is well defined.
\code{\_sample} draws $u$ uniformly and calls \code{\_rescale}, so random draws
and the prior transform share one code path.

\subsection{Is \code{ln\_prob} the density that \code{rescale} samples from?}
\label{sec:discretisation}

Both describe the same KDE, but they discretise it differently.
\code{rescale} inverts piecewise-linear CDFs, so the density it actually
samples is piecewise constant: on each mass cell,
\begin{equation}
  p_{\mathrm{r}}(m) = \frac{F(m_{i+1}) - F(m_i)}{m_{i+1} - m_i},
  \qquad
  p_{\mathrm{r}}(\ell \mid m) = \frac{G(\ell_{j+1} \mid m) - G(\ell_j \mid m)}{\ell_{j+1} - \ell_j},
  \label{eq:rescale_density}
\end{equation}
with $G(\cdot \mid m)$ the blended row of Eq.~\eqref{eq:rowblend}.
\code{ln\_prob} instead interpolates $\ln p$ bilinearly. The two coincide at the
grid points and differ inside cells, most where $\ln p$ changes quickly.

The size of the difference follows from importance-weighting draws from
\code{rescale} by $w = p_{\mathrm{lnprob}} / p_{\mathrm{r}}$, both
evaluated in $(m, \ell)$. Then $\mathbb{E}[w] = \int p_{\mathrm{lnprob}}$ is the
normalisation of $\exp(\code{ln\_prob})$, and the $w$-weighted quantiles are
the quantiles a sampler driven by \code{ln\_prob} (MCMC) converges to. With
$2\times10^6$ draws:
\begin{center}
\begin{tabular}{lcc}
\toprule
& $\int \exp(\code{ln\_prob})$ & largest quantile shift, 5--99\% \\
\midrule
\code{TOVJointDist} $(m, \Lambda)$ & $0.9987$ & $0.3\%$ (in $\Lambda$ at 99\%) \\
\code{TOVMassLambdaTildeJointDist} & $0.9975$ & $0.3\%$ (in $\dlt$ at 99\%) \\
\bottomrule
\end{tabular}
\end{center}
For practical purposes the two are the same density. A nested sampler (driven
by \code{rescale}) and an MCMC sampler (driven by \code{ln\_prob}) target the
same prior to well under a percent, and evidences computed with
$\exp(\code{ln\_prob})$ are normalised to $0.25\%$.

\begin{caveat}{the density is not strictly positive on its own support}
Linear inversion occasionally places a draw in a cell where the KDE was
floored, so $\ln p \approx -690$ there. This affects about $10^{-5}$ of draws
(6 in $4\times10^5$ for the four-parameter prior, with minimum
$\ln p = -626$). It is harmless for sampling, but it makes estimators that
divide by $p$, such as a smooth importance sampler $\mathbb{E}_p[q/p]$ with a
Gaussian $q$, diverge. Weight-based diagnostics need a truncated proposal; the
box estimator of Section~\ref{sec:verify} does not have this problem.
\end{caveat}

\subsection{Using it for a binary}
\label{sec:pairs}

For a binary, create two independent instances, one per star:
\begin{verbatim}
    d1 = TOVJointDist(grid, names=["mass_1", "lambda_1"])
    d2 = TOVJointDist(grid, names=["mass_2", "lambda_2"])
    for d in (d1, d2):
        for name in d.names:
            priors[name] = TOVJointPrior(d, name)
\end{verbatim}
This models the two stars as independent draws from the EOS population, the
same assumption as Eq.~\eqref{eq:component} below. It does \emph{not} impose
$m_1 \ge m_2$: half of all draws have \code{mass\_2 > mass\_1}. Runs that rely
on bilby's labelling convention should add a \code{Constraint} on
\code{mass\_ratio}, or use the four-parameter class, which imposes the ordering
exactly.

The masses here come from the TOV distribution itself. To keep a run's own mass
prior and apply the EOS only to $\Lambda$, use
\code{TOVConditionalLambdaPrior} (Section~\ref{sec:conditional}). Registering
a \code{TOVJointPrior} for \code{lambda\_1} alone, alongside an ordinary
\code{mass\_1} prior, does not do this. Because the joint prior waits for a mass
it never receives, \code{rescale} fails with an inhomogeneous-shape error,
\code{ln\_prob} silently returns $0$ for every $\Lambda$, and \code{sample()}
draws $\Lambda$ at a hidden mass of its own, uncorrelated with \code{mass\_1}.

\begin{bugnote}{sibling instances merged during \code{rescale}}
\code{ConditionalPriorDict.rescale} groups the parameters of a joint
distribution by the string \code{dist.distname}, not by object identity. The
base class sets \code{distname = "joint\_dist"} for every instance, so the
\code{mass\_1}/\code{lambda\_1} and \code{mass\_2}/\code{lambda\_2} pairs would
be merged and their results mixed up. \code{TOVJointDist.\_\_init\_\_} now sets
\code{distname = "tov\_joint\_" + "\_".join(names)}, which is unique per pair.
\code{TOVMassLambdaTildeJointDist} does the same, with prefix
\code{"tov\_mc\_q\_lambda\_tilde\_"}.
\end{bugnote}

%======================================================================
\section{Binary: \code{TOVMassLambdaTildeJointDist}}
\label{sec:binary}
%======================================================================

BNS runs are usually parametrised by $(\Mc, q, \lt, \dlt)$, not by the
component values. Giving $\lt$ and $\dlt$ independent uniform priors throws
away the EOS correlation. This class builds the prior on
$(\Mc, q, \lt, \dlt)$ induced by the same grid. The construction is exact: both
maps involved are invertible, so the density follows by change of variables,
with no approximation beyond the grid itself.

\subsection{Construction}

\code{TOVMassLambdaTildeJointDist.\_\_init\_\_} takes the same \code{grid\_file}
and builds a private \code{TOVJointDist(grid\_file, names=["mass", "lambda"])}
as \code{self.\_grid}. Every grid operation below goes through that object's
methods of Section~\ref{sec:tables}, so the two classes share one
implementation of the tables and of Eq.~\eqref{eq:log10jac}. The internal names
never reach the sampler. The default parameter \code{names} are
\code{["chirp\_mass", "mass\_ratio", "lambda\_tilde", "delta\_lambda\_tilde"]},
in that order; custom names must keep that order.

\subsection{The model in component space}

Write $\theta = (m_1, m_2, \Lambda_1, \Lambda_2)$. The two stars are
independent draws from the same EOS population, conditioned on the labelling
convention $m_1 \ge m_2$:
\begin{equation}
  p(\theta) \;=\; 2\, p(m_1, \Lambda_1)\, p(m_2, \Lambda_2)\,
  \mathbb{1}[m_1 \ge m_2].
  \label{eq:component}
\end{equation}
The factor of $2$ is the usual ordering factor: restricting an exchangeable pair
to the half-space $m_1 \ge m_2$ doubles the density there.

This is a modelling choice worth stating plainly. It treats the two stars as
independent samples from the \emph{population} of EOS-drawn stars, not as two
stars sharing one EOS. A common-EOS prior would correlate $\Lambda_1$ and
$\Lambda_2$ beyond what their masses already imply; Eq.~\eqref{eq:component}
does not.

\subsection{The map to sampling parameters}

Let $x = (\Mc, q, \lt, \dlt)$ with
\begin{equation}
  \Mc = \frac{(m_1 m_2)^{3/5}}{(m_1+m_2)^{1/5}},
  \qquad
  q = \frac{m_2}{m_1} \in (0, 1],
  \qquad
  \eta = \frac{m_1 m_2}{(m_1+m_2)^2} = \frac{q}{(1+q)^2}.
\end{equation}
The tidal combinations are those of Wade et al.\
(\href{https://arxiv.org/abs/1402.5156}{arXiv:1402.5156}), as implemented in
\code{bilby.gw.conversion.lambda\_1\_lambda\_2\_to\_lambda\_tilde} and
\code{lambda\_1\_lambda\_2\_to\_delta\_lambda\_tilde}:
\begin{align}
  \lt  &= \tfrac{8}{13}\Big[ c_1 (\Lambda_1 + \Lambda_2) + c_2 (\Lambda_1 - \Lambda_2) \Big], \\
  \dlt &= \tfrac{1}{2}\Big[ c_3 (\Lambda_1 + \Lambda_2) + c_4 (\Lambda_1 - \Lambda_2) \Big],
\end{align}
with $\eta$-dependent coefficients, returned by
\code{\_lambda\_tilde\_coefficients(eta)}:
\begin{align}
  c_1 &= 1 + 7\eta - 31\eta^2, &
  c_2 &= \sqrt{1-4\eta}\,\big(1 + 9\eta - 11\eta^2\big), \nonumber\\
  c_3 &= \sqrt{1-4\eta}\,\Big(1 - \tfrac{13272}{1319}\eta + \tfrac{8944}{1319}\eta^2\Big), &
  c_4 &= 1 - \tfrac{15910}{1319}\eta + \tfrac{32850}{1319}\eta^2 + \tfrac{3380}{1319}\eta^3 .
  \label{eq:coeffs}
\end{align}
The code clips $1-4\eta$ at zero before the square root, so round-off at
$q = 1$ cannot produce a NaN.

\subsection{The Jacobian factorises}

The map $\theta \mapsto x$ has the block structure
\begin{equation}
  \frac{\partial(\Mc, q, \lt, \dlt)}{\partial(m_1, m_2, \Lambda_1, \Lambda_2)}
  =
  \begin{pmatrix}
    \dfrac{\partial(\Mc, q)}{\partial(m_1, m_2)} & \mathbf{0} \\[2ex]
    \dfrac{\partial(\lt, \dlt)}{\partial(m_1, m_2)} &
    \dfrac{\partial(\lt, \dlt)}{\partial(\Lambda_1, \Lambda_2)}
  \end{pmatrix},
\end{equation}
because the masses do not depend on the tidal parameters. The upper-right block
vanishes, so the matrix is block triangular and its determinant is the product
of the two diagonal blocks. The lower-left block, which is nonzero and depends
on $\eta$, drops out entirely.

\subsubsection{Mass block}

With $m_1 = \Mc (1+q)^{1/5} q^{-3/5}$ and $m_2 = q m_1$,
\begin{equation}
  \frac{\partial m_1}{\partial \Mc} = \frac{m_1}{\Mc},
  \qquad
  \frac{\partial m_2}{\partial \Mc} = \frac{m_2}{\Mc} = \frac{q m_1}{\Mc},
  \qquad
  \frac{\partial m_2}{\partial q} = m_1 + q \frac{\partial m_1}{\partial q},
\end{equation}
so
\begin{equation}
  \det \frac{\partial(m_1, m_2)}{\partial(\Mc, q)}
  = \frac{m_1}{\Mc}\left( m_1 + q \frac{\partial m_1}{\partial q} \right)
    - \frac{\partial m_1}{\partial q}\, \frac{q m_1}{\Mc}
  = \frac{m_1^2}{\Mc}.
  \label{eq:massjac}
\end{equation}
The $\partial m_1/\partial q$ terms cancel, leaving the same $m_1^2/\Mc$ that
appears in \code{bilby.gw.prior.convert\_to\_flat\_in\_component\_mass\_prior}.

\subsubsection{Tidal block}

At fixed masses the tidal map is \emph{linear}:
\begin{equation}
  \begin{pmatrix} \lt \\ \dlt \end{pmatrix}
  = A(\eta)
  \begin{pmatrix} \Lambda_1 \\ \Lambda_2 \end{pmatrix},
  \qquad
  A(\eta) =
  \begin{pmatrix}
    \tfrac{8}{13}(c_1 + c_2) & \tfrac{8}{13}(c_1 - c_2) \\[1ex]
    \tfrac{1}{2}(c_3 + c_4)  & \tfrac{1}{2}(c_3 - c_4)
  \end{pmatrix}.
\end{equation}
Hence
\begin{align}
  \det A
  &= \tfrac{8}{13}\cdot\tfrac{1}{2}\Big[ (c_1+c_2)(c_3-c_4) - (c_1-c_2)(c_3+c_4) \Big] \nonumber \\
  &= \tfrac{4}{13}\Big[ -2 c_1 c_4 + 2 c_2 c_3 \Big]
   = \frac{8}{13}\big( c_2 c_3 - c_1 c_4 \big),
  \label{eq:tidaljac}
\end{align}
a function of $\eta$ alone, computed in log form by
\code{\_ln\_abs\_det\_lambda\_tilde\_jacobian(eta)}. At equal masses
$\eta \to 1/4$ and $\sqrt{1-4\eta} \to 0$, so $c_2, c_3 \to 0$ and
$\det A \to -\tfrac{8}{13} c_1 c_4 \neq 0$. The map stays invertible on the whole
physical range, and no special case is needed at $q = 1$.

\subsection{The density}
\label{sec:binary_density}

Combining Eqs.~\eqref{eq:component}, \eqref{eq:massjac} and \eqref{eq:tidaljac},
\begin{equation}
  \boxed{\;
  p(\Mc, q, \lt, \dlt)
  = 2\, p(m_1, \Lambda_1)\, p(m_2, \Lambda_2)\,
    \frac{m_1^2}{\Mc}\, \frac{1}{\lvert \det A(\eta) \rvert}
  \;}
\end{equation}
or, term by term as \code{TOVMassLambdaTildeJointDist.\_ln\_prob} evaluates it,
\begin{align}
  \ln p(\Mc, q, \lt, \dlt) =\;&
    \ln 2
    + \ln p(m_1, \Lambda_1) + \ln p(m_2, \Lambda_2) \nonumber \\
    &+ 2\ln m_1 - \ln \Mc
    - \ln\big\lvert \det A(\eta) \big\rvert .
  \label{eq:binary_lnprob}
\end{align}
The component values are recovered from $x$ with the standard inversions,
\code{chirp\_mass\_and\_mass\_ratio\_to\_total\_mass},
\code{total\_mass\_and\_mass\_ratio\_to\_component\_masses} and
\code{lambda\_tilde\_delta\_lambda\_tilde\_to\_lambda\_1\_lambda\_2} (all in
\code{bilby.gw.conversion}). Each $\ln p(m_k, \Lambda_k)$ comes from
\code{self.\_grid.ln\_pdf\_mass\_lambda}, so it already includes the corrected
Jacobian of Eq.~\eqref{eq:log10jac}. The bug of Section~\ref{sec:log10jac}
entered this formula twice, once per star.

\subsection{Support and bounds}
\label{sec:support}

The prior is \emph{not} supported on a box in $x$. A rectangle in
$(\lt, \dlt)$ contains points whose implied $\Lambda_1$ or $\Lambda_2$ is
negative, and points that fall outside the EOS grid entirely.
\code{\_ln\_prob} builds a \code{valid} mask that requires $q \in (0, 1]$,
$\Mc > 0$, $\Lambda_1, \Lambda_2 > 0$, and finite recovered components. At
invalid points it substitutes harmless in-grid values before calling the
interpolant, which keeps NaNs out of it, and then sets those points to
$-\infty$. Points that are valid but off the grid get $-\infty$ from the
interpolant's \code{fill\_value}.

\code{TOVMassLambdaTildeJointDist.default\_bounds()} returns \emph{outer}
bounds on this non-rectangular support; being inside them does not imply
finite probability. They are derived as follows.
\begin{itemize}
  \item $\Mc$ increases in both masses, so it is bracketed by
    $\Mc(m_{\min}, m_{\min})$ and $\Mc(m_{\max}, m_{\max})$.
  \item $q \in [m_{\min}/m_{\max},\, 1]$.
  \item At fixed masses, $\lt$ and $\dlt$ are linear in $(\Lambda_1, \Lambda_2)$,
    so their extrema over the grid's $\Lambda$ range lie at the four corners
    $\Lambda_k \in \{10^{\ell_{\min}}, 10^{\ell_{\max}}\}$. These are scanned
    over a $200 \times 200$ grid of ordered masses $m_1 \ge m_2$.
\end{itemize}
The tidal bounds come from a scan, not a proof. Empirically, none of
$4\times10^5$ draws from \code{rescale} fell outside them.

\begin{caveat}{exact unit-cube edges}
At $u_k \in \{0, 1\}$ exactly, \code{rescale} returns a point on the grid
boundary. The round trip through the inversions above can then land
$\sim10^{-16}$ outside it, e.g.\ $m_1 = m_{\min} - 1.1\times10^{-16}$, and
\code{ln\_prob} returns $-\infty$. Continuous uniform draws essentially never hit
these values; the interior of the cube is unaffected.
\end{caveat}

\subsection{Prior transform}
\label{sec:binary_rescale}

\code{TOVMassLambdaTildeJointDist.\_component\_rescale} maps
$u \in [0,1]^4$ to $\theta$, and \code{\_rescale} then applies
$\theta \mapsto x$ with \code{component\_masses\_to\_chirp\_mass},
$q = m_2/m_1$ and the two tidal conversions.

The ordering $m_1 \ge m_2$ is imposed by the transform itself, not by rejection,
which keeps the map a smooth bijection. If $m_1$ is the larger of two
independent draws from $p(m)$, its CDF is $F(m)^2$, so
\begin{equation}
  m_1 = F^{-1}\!\left( \sqrt{u_0} \right),
  \qquad
  m_2 = F^{-1}\!\big( u_1\, F(m_1) \big),
  \label{eq:ordered}
\end{equation}
the second being $p(m)$ truncated to $[m_{\min}, m_1]$. In code, these are
\code{\_grid.mass\_from\_cdf(np.sqrt(u0))} and
\code{\_grid.mass\_from\_cdf(u1 * \_grid.mass\_cdf(m1))}. The induced density is
\begin{equation}
  \underbrace{2\, p(m_1) F(m_1)}_{\text{max of two draws}} \times
  \underbrace{\frac{p(m_2)}{F(m_1)}}_{\text{truncated}}
  = 2\, p(m_1)\, p(m_2),
\end{equation}
exactly the mass part of Eq.~\eqref{eq:component}; the $F(m_1)$ factors cancel.
They also cancel for the discretised $F$ the code actually uses, so the
ordering adds no discretisation error beyond that of
Section~\ref{sec:discretisation}. The deformabilities then follow
star by star from Eq.~\eqref{eq:single_rescale},
\begin{equation}
  \Lambda_1 = 10^{\,G^{-1}(u_2 \mid m_1)},
  \qquad
  \Lambda_2 = 10^{\,G^{-1}(u_3 \mid m_2)},
\end{equation}
via \code{\_grid.log10\_lambda\_from\_cdf}.

Which cube coordinate feeds which physical parameter is arbitrary.
\code{rescale} maps the whole vector at once, so the sampler never relies on a
per-coordinate correspondence. As with the single-star class, \code{\_sample}
is \code{\_rescale} applied to uniform draws.

%======================================================================
\section{Conditional deformability: \code{TOVConditionalLambdaPrior}}
\label{sec:conditional}
%======================================================================

Both joint classes also fix the mass prior, to the EOS population's mass
distribution. A run often has its own mass prior, e.g.\ uniform in component
masses, and wants the EOS only to shape $\Lambda$ given that mass.
\code{TOVConditionalLambdaPrior} does this.

\subsection{Model}

With any mass prior $\pi(m_1, m_2)$ chosen by the user, the prior is
\begin{equation}
  p(m_1, m_2, \Lambda_1, \Lambda_2)
  = \pi(m_1, m_2)\; p(\Lambda_1 \mid m_1)\; p(\Lambda_2 \mid m_2),
  \qquad
  p(\Lambda \mid m) = \frac{p_{\mathrm{TOV}}(m, \Lambda)}{p_{\mathrm{TOV}}(m)} .
  \label{eq:conditional}
\end{equation}
One instance is created per star (\code{component=1} or \code{2}, default name
\code{lambda\_<component>}). Each depends only on its own star's mass, so, as in
Eq.~\eqref{eq:component}, the two deformabilities are independent given the
masses. Because the masses are conditioned on rather than transformed, no mass
Jacobian enters. The prior on the masses is exactly $\pi$, and $\Lambda$ is
correlated with the mass the run actually samples.

\subsection{The conditional density, exactly normalised}

The joint density $p_{\mathrm{TOV}}(m, \Lambda)$ is the one every other class
uses, \code{TOVJointDist.ln\_pdf\_mass\_lambda} (Eq.~\ref{eq:single}). Its
normaliser $p_{\mathrm{TOV}}(m) = \int p_{\mathrm{TOV}}(m, \Lambda)\,\dd\Lambda$
is computed from the same interpolant, not from the trapezoid table of
Eq.~\eqref{eq:massmarg}, so that each conditional integrates to exactly one.

The interpolant is bilinear in $\ln p$. At fixed mass $m$ in the cell
$[m_i, m_{i+1}]$, it is therefore piecewise \emph{linear} in $\ell$, with node
values blended linearly in mass:
\begin{equation}
  \ln p_{\mathrm{TOV}}(m, \ell) = a_j + d_j\, t,
  \qquad
  a_j = (1-w)\,\code{log\_pdf}[i, j] + w\,\code{log\_pdf}[i+1, j],
  \qquad
  d_j = a_{j+1} - a_j,
  \label{eq:rowlinear}
\end{equation}
for $\ell = \ell_j + t\,h_j$ with $t \in [0, 1]$, cell width
$h_j = \ell_{j+1} - \ell_j$, and $w$ as in Eq.~\eqref{eq:rowblend}. The density
is piecewise exponential, so each cell integrates in closed form:
\begin{equation}
  I_j = \int_{\ell_j}^{\ell_{j+1}} p_{\mathrm{TOV}}(m, \ell)\,\dd\ell
      = h_j\,\frac{e^{a_{j+1}} - e^{a_j}}{d_j}
  \quad\big(\to h_j\, e^{a_j} \text{ as } d_j \to 0\big),
  \qquad
  p_{\mathrm{TOV}}(m) = \sum_j I_j .
  \label{eq:cellint}
\end{equation}
Then $\int p_{\mathrm{TOV}}(m,\Lambda)\,\dd\Lambda = \int p_{\mathrm{TOV}}(m,\ell)\,\dd\ell$,
because the $\ell \to \Lambda$ Jacobian of Eq.~\eqref{eq:log10jac} cancels.
\code{ln\_prob} returns
\begin{equation}
  \ln p(\Lambda \mid m) = \ln p_{\mathrm{TOV}}(m, \Lambda) - \ln \sum_j I_j(m) .
\end{equation}
Numerically, each row is shifted by its maximum before exponentiating. The
grid's log-density spans $[-690.8, 0.33]$, with jumps of up to $685$ between
neighbouring cells, so the unshifted form would overflow or underflow.
\code{TOVConditionalLambdaPrior.\_conditional\_tables} computes the shifted
$a_j$, the cumulative sums $C_j = \sum_{k<j} I_k$, and $\ln p_{\mathrm{TOV}}(m)$,
vectorised over masses.

\subsection{Prior transform and CDF}

The same closed form gives the CDF and its inverse, with no separate tables.
For $\ell$ in cell $j$ at fraction $t$,
\begin{equation}
  G(\ell \mid m) = \frac{C_j + h_j\,\big(e^{a_j + d_j t} - e^{a_j}\big)/d_j}{\sum_k I_k},
  \label{eq:condcdf_exact}
\end{equation}
implemented by \code{TOVConditionalLambdaPrior.cdf}. For \code{rescale}, the
target $r = u \sum_k I_k$ selects the cell with $C_j \le r < C_{j+1}$. The
remainder $\rho = r - C_j$ then gives
\begin{equation}
  t = \frac{\ln\!\big(e^{a_j} + \rho\, d_j / h_j\big) - a_j}{d_j}
  \quad\big(\to \rho / (h_j e^{a_j}) \text{ as } d_j \to 0\big),
  \qquad
  \Lambda = 10^{\ell_j + t h_j} .
  \label{eq:condinv}
\end{equation}
The argument of the logarithm is evaluated directly, without dividing by
$e^{a_j}$, so cells at the KDE floor are handled without overflow. Cells of
zero width in probability are never selected.

\paragraph{Consequence.}
Unlike the joint classes (Section~\ref{sec:discretisation}), \code{rescale} and
\code{ln\_prob} here describe \emph{the same} density, not two
discretisations of it. The verification in Section~\ref{sec:verify_conditional}
confirms this to rounding error.

\subsection{Which parameters it conditions on}

\code{mass\_parameters} selects the required variables, i.e.\ the keys
\code{ConditionalPriorDict} must supply:
\begin{itemize}
  \item \code{"component\_masses"} (default): \code{mass\_<component>}, for runs
    that sample the component masses directly.
  \item \code{"chirp\_mass\_and\_mass\_ratio"}: \code{chirp\_mass} and
    \code{mass\_ratio}, converted internally with
    \code{chirp\_mass\_and\_mass\_ratio\_to\_total\_mass} and
    \code{total\_mass\_and\_mass\_ratio\_to\_component\_masses}. This mode is
    for runs such as \code{prior\_files/aligned\_spins\_bns.prior}, which
    sample $(\Mc, q)$ and hold \code{mass\_1}, \code{mass\_2} as
    \code{Constraint}s. A constraint is not a sampled parameter, so it cannot
    be conditioned on, and \code{ConditionalPriorDict} would report the
    conditions as unresolvable. As the conditioning is on a deterministic
    function of $(\Mc, q)$, no Jacobian is needed here either.
\end{itemize}
For a given mass, both modes return the same density.

\subsection{Masses outside the grid}

$p(\Lambda \mid m)$ is defined only for $m \in [1.000, 3.476]\,M_\odot$.
Outside that range, \code{ln\_prob} returns $-\infty$, and \code{rescale}
returns \code{nan} with a one-off warning. It does not raise, because bilby
rescales every key before checking constraints, and an out-of-grid draw that a
\code{Constraint} is about to reject must not abort the run. The mass prior
should therefore lie inside the grid. In
\code{"chirp\_mass\_and\_mass\_ratio"} mode, the component masses implied by a
$(\Mc, q)$ box generally reach beyond it. The \code{mass\_1} and \code{mass\_2}
constraints must then be tightened to the grid, e.g.\ to $[1.0, 3.47]$,
instead of the $[0.5, 5]$ in \code{aligned\_spins\_bns.prior}.

\subsection{Implementation notes}

\begin{itemize}
  \item The class follows the idiom of \code{ConditionalChiInPlane} in the same
    file. It subclasses \code{ConditionalBasePrior}, sets
    \code{\_required\_variables} explicitly, passes a no-op
    \code{\_condition\_function}, and overrides \code{rescale},
    \code{ln\_prob}, \code{prob} and \code{cdf}. \code{sample} is inherited
    and calls \code{rescale} on uniform draws. Calls with the wrong required
    variables raise \code{IllegalRequiredVariablesException}. A call with none,
    i.e.\ outside a \code{ConditionalPriorDict}, raises a \code{ValueError}
    saying so.
  \item \code{conditional\_prior\_factory} renames the class on instantiation.
    The constructor restores the name, and \code{\_\_repr\_\_} and
    \code{get\_instantiation\_dict} use the \code{Prior} versions, so the prior
    round-trips through prior files and result JSON. \code{grid\_file} is stored
    as given; use an absolute path if results will be reloaded elsewhere.
  \item \code{grid\_file} may also be an existing \code{TOVJointDist}, whose
    loaded grid is then shared. \code{TOVJointDist.\_\_init\_\_} keeps the raw
    table as \code{\_log\_pdf} for this class.
  \item Scalar inputs return NumPy scalars, not 0-d arrays, for the LAL
    \code{REAL8} reason given in Section~\ref{sec:machinery}.
\end{itemize}

%======================================================================
\section{Integration with bilby's joint-prior machinery}
\label{sec:machinery}
%======================================================================

\subsection{How a sampler reaches the classes}

Each parameter has its own \code{Prior} object (\code{TOVJointPrior} or
\code{TOVMassLambdaTildeJointPrior}), all sharing one \code{Dist}. Their
constructors only add a type check on \code{dist} to
\code{JointPrior.\_\_init\_\_}. The shared \code{Dist} holds per-call state,
and the inherited \code{JointPrior} methods use it as follows.
\begin{itemize}
  \item \code{rescale(val)} stores \code{val} in \code{dist.rescale\_parameters}
    and returns \code{[]} until every parameter of the dist has been
    supplied. It then stacks the stored values \emph{in \code{dist.names}
    order}, calls \code{dist.rescale}, and resets the store. The order in
    which a sampler lists its search keys therefore does not matter.
    \code{ConditionalPriorDict.rescale} then copies the full vector back to
    each parameter, grouping by \code{distname} (Section~\ref{sec:pairs}).
  \item \code{ln\_prob(val)} works the same way through
    \code{dist.requested\_parameters}. Every parameter but the last returns $0$,
    and the last returns the joint log-density, so the sum
    \code{PriorDict.ln\_prob} forms is correct.
  \item \code{sample()} draws all parameters of the dist at once on the first
    call and hands out the stored values on later calls, so the parameters of
    one draw belong to the same $(m, \Lambda)$ or $\theta$.
\end{itemize}
Nested samplers (\code{dynesty} and others) use only \code{rescale}; MCMC samplers
(\code{emcee}, \code{bilby\_mcmc}) use only \code{ln\_prob}. Both feed the
\code{log\_prior} column of the result, which \code{Result} computes as
\code{priors.ln\_prob(dict(posterior), axis=0)}, i.e.\ vectorised over samples.

\paragraph{The conditional prior takes a different route.}
\code{TOVConditionalLambdaPrior} holds no shared state and goes through
\code{ConditionalPriorDict}'s conditional path. The dict recognises it by its
\code{condition\_func} attribute and orders it after the keys listed in
\code{required\_variables}. It then passes their most recently sampled values
(\code{least\_recently\_sampled}) as keyword arguments to \code{rescale},
\code{ln\_prob} and \code{sample}. For vectorised \code{ln\_prob}, these are
arrays, and the class broadcasts them against the $\Lambda$ values.

\subsection{Fixes to the interface}

Four interface bugs made the original classes fail inside real runs. All are
fixed in the subclasses, leaving \code{bilby/core/prior/joint.py} untouched so
that the fork stays a clean addition on top of upstream. Each fix is harmless if
upstream changes later.

\begin{bugnote}{\code{xp} keyword}
Upstream \code{BaseJointPriorDist.ln\_prob} now passes an array-namespace
argument \code{xp} to \code{\_ln\_prob}. Both \code{\_ln\_prob} implementations
are decorated with \code{@xp\_wrap} and accept \code{xp}. They compute in NumPy,
because \code{RegularGridInterpolator} is NumPy-only, and convert back with
\code{xp.asarray}. The classes are therefore NumPy-bound: an \code{xp} of CuPy
or JAX fails in \code{\_ln\_prob}.
\end{bugnote}

\begin{bugnote}{\code{random\_state} ignored}
\code{\_sample} drew from NumPy's global generator, so seeded runs did not
reproduce and pool workers could share a stream. Both \code{\_sample}
methods now resolve the generator with
\code{bilby.core.utils.random.resolve\_random\_state(random\_state)}.
\end{bugnote}

\begin{bugnote}{0-d arrays from \code{sample()}}
\code{JointPrior.sample} ends in \code{sample.squeeze()}, which for a size-1
draw returns a 0-d \code{ndarray}, not a float. LAL's SWIG \code{REAL8} typemap
rejects it:
\begin{center}
\small\code{TypeError: in method}\\\code{'SimInspiralWaveformParamsInsertTidalLambda1',}\\\code{argument 2 of type 'REAL8'}
\end{center}
The masses escaped only because they are multiplied by \code{lal.MSUN\_SI}
before reaching LAL, which turns them into scalars. \code{TOVJointPrior.sample}
and \code{TOVMassLambdaTildeJointPrior.sample} append \code{[()]}, which unwraps
0-d arrays and leaves larger results unchanged.
\end{bugnote}

\begin{bugnote}{shape-$(1,)$ arrays from \code{ln\_prob}}
\code{BaseJointPriorDist.ln\_prob} returns shape $(1,)$ for a single sample,
while ordinary priors return floats. \code{PriorDict.ln\_prob} then fails while
stacking the per-parameter terms:
\begin{center}
\code{ValueError: setting an array element with a sequence. ...}
\end{center}
\code{TOVJointDist.ln\_prob} and \code{TOVMassLambdaTildeJointDist.ln\_prob}
override it to return a scalar in that case, matching what
\code{BaseJointPriorDist.rescale} already does.
\end{bugnote}

\subsection{Bounds}

\begin{caveat}{tightened bounds are honoured by \code{ln\_prob} only}
Setting a narrower range on one parameter, e.g.\
\code{priors["lambda\_tilde"].maximum = 600}, updates \code{dist.bounds}.
\code{BaseJointPriorDist.\_check\_samp} then sets \code{ln\_prob} to $-\infty$
outside the range, but \code{BaseJointPriorDist.rescale} applies no bounds. In
that example, $26\%$ of \code{rescale} draws still have $\lt > 600$, so a nested
sampler and an MCMC sampler would explore different priors. The truncated
density is also no longer normalised.

There is a second, upstream problem in \code{\_check\_samp}: for a batch of
samples, it keeps the out-of-bounds mask of the \emph{first} parameter that has
any violation and does not combine masks across parameters. With bounds
tightened on two parameters, the vectorised \code{ln\_prob} used for
\code{log\_prior} let 349 of 532 violating points through with finite
probability. Evaluated one sample at a time, all of them correctly get
$-\infty$.

Neither problem arises with the default bounds, where the support check in
\code{\_ln\_prob} does the real work. To restrict the range, leave the bounds
alone and add a \code{bilby.core.prior.Constraint} instead. Constraints are
applied consistently and renormalised by \code{PriorDict}.
\end{caveat}

%======================================================================
\section{Verification}
\label{sec:verify}
%======================================================================

\subsection{\code{verify\_tov\_joint\_priors.py}}

The script at the repository root runs four groups of checks on the
\code{Dist} classes directly:
\begin{verbatim}
    python verify_tov_joint_priors.py [path/to/tov_mass_lambda_joint_kde_grid.npz]
\end{verbatim}

\paragraph{1. Jacobians against finite differences.}
At four configurations, including near-equal masses where
$\sqrt{1-4\eta} \to 0$, the full $4\times4$ determinant of $\theta \mapsto x$ is
compared with the analytic $m_1^2/\Mc \cdot \lvert\det A\rvert^{-1}$. They agree
to better than $10^{-6}$ in the log, which also confirms the block-triangularity
argument.

\paragraph{2. Ordering and invertibility.}
For $2\times10^4$ draws from \code{\_component\_rescale}, $m_1 \ge m_2$ and
$q \in (0, 1]$ always hold. Pushing the draws back through the inversions of
\code{\_ln\_prob} recovers $\theta$ to a relative error of $1.5\times10^{-11}$,
and \code{ln\_prob} is finite on all of them.

\paragraph{3. \code{ln\_prob} is the density \code{rescale} samples from.}
This is the check that caught the bug of Section~\ref{sec:log10jac}. For
samples $x_i \sim p$ and any box $B$,
\begin{equation}
  \frac{1}{N} \sum_{x_i \in B} \frac{1}{\mathrm{vol}(B)\, p(x_i)}
  \;\xrightarrow{\;N \to \infty\;}\;
  \int_B \frac{p(x)}{\mathrm{vol}(B)\, p(x)}\,\dd x \;=\; 1 .
\end{equation}
This unbiased estimator of $1$ does not depend on how sharply $p$ varies inside
$B$. A raw histogram comparison does, and for this strongly correlated density
it is off by more than an order of magnitude at feasible sample sizes. Any
missing constant, Jacobian or normalisation shows up directly as a departure
from $1$.

\paragraph{4. Control.}
The same estimator applied to the two-parameter \code{TOVJointDist}.

\begin{table}[h]
\centering
\begin{tabular}{lrr}
\toprule
Box centre & before fix & current code \\
\midrule
\multicolumn{3}{l}{\code{TOVMassLambdaTildeJointDist}, $(\Mc,\, q,\, \lt,\, \dlt)$} \\
$(1.20,\ 0.85,\ 400,\ -30)$ & $17.9 \pm 1.2$ & $1.02 \pm 0.07$ \\
$(1.35,\ 0.75,\ 250,\ -40)$ & $19.0 \pm 3.1$ & $0.96 \pm 0.08$ \\
$(1.15,\ 0.90,\ 600,\ -20)$ & $18.8 \pm 1.2$ & $1.06 \pm 0.07$ \\
\midrule
\multicolumn{3}{l}{control: \code{TOVJointDist}, $(m, \Lambda)$} \\
$(1.30,\ 400)$ & --- & $1.02 \pm 0.02$ \\
$(1.80,\ 100)$ & --- & $1.00 \pm 0.01$ \\
\bottomrule
\end{tabular}
\caption{Box estimator, target $1$, with $N = 4\times10^5$ samples. The
``before fix'' column was recorded when the bug was found. The constant
$18.8$ is $(10/\ln 10)^2$, the same error once per star.}
\end{table}

The box estimator only probes the bulk of the distribution. The global
comparison of Section~\ref{sec:discretisation}, done by importance-weighting
draws with the exact density of Eq.~\eqref{eq:rescale_density}, extends the
agreement to the whole support, tails included.

\subsection{Checks through bilby's sampling interface}

The script calls the \code{Dist} classes directly. The following was also
checked through \code{BNSPriorDict} and \code{bilby.run\_sampler}, exactly as a
run calls the classes.
\begin{itemize}
  \item \code{ConditionalPriorDict.rescale} reproduces \code{dist.rescale} for
    any ordering of the search keys, with fixed and ordinary priors mixed in,
    and leaves no state behind. \code{PriorDict.ln\_prob} returns a scalar for
    single samples and the correct array for the vectorised
    \code{log\_prior} computation. Single draws are scalars, and the four
    values of one draw are coherent.
  \item \textbf{dynesty}, four-parameter prior with a Gaussian likelihood on
    $(\Mc, \lt)$: $\ln Z = -1.579 \pm 0.043$ against a Monte Carlo value of
    $-1.562$. All 5/50/95\% posterior quantiles of the four parameters match
    the importance-reweighted prior, and the \code{log\_prior} column equals
    \code{dist.ln\_prob}. The same holds with \code{npool=2}
    ($\ln Z = -1.583 \pm 0.059$), which confirms the classes pickle and work
    in spawned workers.
  \item \textbf{dynesty}, two \code{TOVJointDist} pairs:
    $\ln Z = -1.743 \pm 0.063$ against $-1.734$, with quantiles in agreement.
  \item \textbf{emcee} with a flat likelihood reproduces the prior's $\Mc$ and
    $q$ to $2\%$ of the 90\% width. It shows tails in $\lt$ and $\dlt$ that are
    $2$--$8\%$ heavy, which is slow mixing of the stretch move in long,
    correlated tails: Section~\ref{sec:discretisation} bounds any genuine
    \code{ln\_prob}/\code{rescale} difference at $0.3\%$. MCMC runs dominated by
    the prior need long, thinned chains.
\end{itemize}
\subsection{\code{TOVConditionalLambdaPrior}}
\label{sec:verify_conditional}

\paragraph{Exactness.}
\begin{itemize}
  \item $\int \exp(\code{ln\_prob})\,\dd\Lambda = 1$ to $4.5\times10^{-8}$ at eight
    masses, including both grid edges. The residual comes from the fine
    numerical quadrature used to check it.
  \item $\lvert \code{cdf}(\code{rescale}(u)) - u \rvert < 1.3\times10^{-12}$
    over $2\times10^5$ random $(u, m)$.
  \item \code{cdf} agrees with a direct numerical integral of
    $\exp(\code{ln\_prob})$ to $3\times10^{-8}$. Together with the previous
    item, this shows \code{rescale} samples exactly the density \code{ln\_prob}
    returns. KS tests of \code{rescale} draws against \code{cdf} at three masses
    give $p = 0.43$--$0.81$.
  \item $\ln p_{\mathrm{joint}}(m, \Lambda) - \ln p(\Lambda \mid m)$, using
    \code{TOVJointDist.ln\_prob} for the joint, is constant in $\Lambda$ to
    $4\times10^{-14}$. The conditional is therefore the joint renormalised, as
    Eq.~\eqref{eq:conditional} states.
\end{itemize}

\paragraph{Through \code{BNSPriorDict}.}
\begin{itemize}
  \item The $\Lambda$ keys resolve as conditional.
  \item \code{rescale} and \code{ln\_prob} match direct calls for any key order.
  \item \code{ln\_prob} is a scalar for single samples and correct when
    vectorised.
  \item Single draws are NumPy scalars.
  \item \code{mass\_1} keeps its own \code{Uniform} prior (KS), while
    $\mathrm{corr}(m_1, \ln\Lambda_1) = -0.89$, so $\Lambda$ follows the run's
    mass.
  \item The \code{"chirp\_mass\_and\_mass\_ratio"} mode gives the same density
    as the component mode at the implied mass.
  \item Off-grid masses give \code{nan} with a warning (\code{rescale}) and
    $-\infty$ (\code{ln\_prob}).
  \item Wrong or missing required variables raise.
  \item The prior survives repr, a prior-file round trip and a JSON round trip.
\end{itemize}

\paragraph{Samplers.}
With uniform component masses on $[1, 2]$ and a Gaussian likelihood on
$(\Lambda_1, \Lambda_2)$:
\begin{center}
\begin{tabular}{lcc}
\toprule
mass parameters & dynesty $\ln Z$ & Monte Carlo $\ln Z$ \\
\midrule
\code{component\_masses} & $-1.014 \pm 0.034$ & $-1.003$ \\
\code{chirp\_mass\_and\_mass\_ratio} & $-1.166 \pm 0.035$ & $-1.172$ \\
\bottomrule
\end{tabular}
\end{center}
In both modes, all 5/50/95\% posterior quantiles of the masses and
deformabilities match the importance-reweighted prior. The \code{log\_prior}
column is finite, and the result reloads from JSON with the prior intact.
\textbf{emcee} with a flat likelihood, which uses only \code{ln\_prob},
reproduces the prior. With 64 walkers and $3\times10^4$ steps, the
5--95\% quantiles deviate by $0.4\%$ of the 90\% width for $m_1$ and $2.4\%$ for
$\Lambda_1$. A $5\times$ shorter chain gave $5.1\%$ for $\Lambda_1$, with the
upper tail on the other side of the reference. The residual is mixing in the
heavy $\Lambda$ tail, not a density error, which the exactness checks above
rule out.

The fork's \code{requirements.txt} requires \code{dynesty>=3.0.0}, because
\code{bilby.core.sampler.dynesty\_utils} imports \code{dynesty.internal\_samplers};
dynesty 2.x fails at import.

%======================================================================
\section{Usage}
%======================================================================

\paragraph{Four-parameter prior,} for runs in $(\Mc, q, \lt, \dlt)$:
\begin{verbatim}
    from bilby.gw.prior import (
        BNSPriorDict, TOVMassLambdaTildeJointDist, TOVMassLambdaTildeJointPrior)

    grid = "prob_data/tov_mass_lambda_joint_kde_grid.npz"
    dist = TOVMassLambdaTildeJointDist(grid)
    priors = BNSPriorDict()   # then add the remaining BNS parameters
    for name in dist.names:
        priors[name] = TOVMassLambdaTildeJointPrior(dist, name)
\end{verbatim}

\paragraph{Per-star joint prior,} for runs in component masses and deformabilities:
see Section~\ref{sec:pairs}, and note that it does not order the masses.

\paragraph{Conditional $\Lambda$ with your own mass prior,} sampling component
masses:
\begin{verbatim}
    from bilby.core.prior import Uniform
    from bilby.gw.prior import BNSPriorDict, TOVConditionalLambdaPrior

    grid = "/abs/path/prob_data/tov_mass_lambda_joint_kde_grid.npz"
    priors = BNSPriorDict(dictionary=dict(
        mass_1=Uniform(1.0, 2.0, "mass_1"),
        mass_2=Uniform(1.0, 2.0, "mass_2"),
        lambda_1=TOVConditionalLambdaPrior(grid, component=1),
        lambda_2=TOVConditionalLambdaPrior(grid, component=2),
    ))  # then add the remaining BNS parameters
\end{verbatim}
or, in a prior file that samples $(\Mc, q)$, as \code{aligned\_spins\_bns.prior}
does:
\begin{verbatim}
    mass_1 = Constraint(name='mass_1', minimum=1.0, maximum=3.47)
    mass_2 = Constraint(name='mass_2', minimum=1.0, maximum=3.47)
    lambda_1 = bilby.gw.prior.TOVConditionalLambdaPrior(grid_file='/path/grid.npz',
        component=1, mass_parameters='chirp_mass_and_mass_ratio')
    lambda_2 = bilby.gw.prior.TOVConditionalLambdaPrior(grid_file='/path/grid.npz',
        component=2, mass_parameters='chirp_mass_and_mass_ratio')
\end{verbatim}

\paragraph{In all cases:}
\begin{itemize}
  \item for the joint classes, use one \code{Dist} instance per joint group and
    pass the same object to each of its \code{Prior}s;
  \item restrict ranges with a \code{Constraint}, not by changing
    \code{minimum}/\code{maximum} (Section~\ref{sec:machinery});
  \item keep every mass the EOS is evaluated at inside the grid,
    $[1.000, 3.476]\,M_\odot$;
  \item the classes are NumPy-only.
\end{itemize}

\end{document}
